\documentclass[11pt]{article}
\usepackage{fiche}
\matiere{Support Vecteur Machine}
\chapitre{M2}
\source{SVM.pdf}
%% statut: propre

\begin{document}
%% page: SVM.pdf
\titre{\couleur{rose}{\entoure{1} \souligne{SVM}}}

\begin{itemize}
\item Données qui appartient à $\mathbb{R}^D$

\couleur{rouge}{\# On veut trouver une frontière \souligne{linéaire} pour séparer les classe}

\item \souligne{Théorème de Cover}

Si on trouve pas de frontière linéaire en $\mathbb{R}^D$ on peut passer en $\mathbb{R}^{D+1}$ pour la trouver

\item L'objectif c'est de trouver un classifieur g selon nos données générées de la sorte :
\[ Y = \entoure{$g$}(X) \]
\couleur{rouge}{g = \souligne{inconnu}, on \entoure{doit} le trouver}

Le classifieur \entoure{g} : $g : X \to \{-1, 1\}$

\couleur{rouge}{\# \entoure{$-1$} = Mal classé \# \entoure{$1$} = Bien classé}
\end{itemize}

\begin{exemple}
\couleur{vert}{$X \in \mathbb{R}^3$, $w = (3, 2, 7)$, $b = -2$ \\
$X_A = (0, 2, -4)$, $Y_A = \entoure{1}$ \\
$h(X_A) = \langle w, x_A \rangle + b = -26$ \\
$g(X_A) = \text{signe}(h(X_A)) = \entoure{$-1$}$ \\
$\hookrightarrow$ prédict° \souligne{mauvaise}}
\end{exemple}

\entoure{g} doit minimiser $\Pr(g(X) \neq Y)$

\entoure{g} est défini comme étant :
\begin{encadre}
\[ g(X) = \text{signe}(h(X)) = \begin{cases} -1 & \text{si } h(x) < 0 \\ 1 & \text{si } h(x) > 0 \end{cases} \]
\end{encadre}

\couleur{rouge}{\# Rappel : $h(X)$ = linéaire donc de la forme
\[ h(X) = \langle w, x \rangle + b = \sum_{i=1}^{p} (w_i x_i) + b \]}

\marge{\couleur{rose}{\souligne{Classifieur}} \\ $g$ minimise $\Pr(g(X) \neq Y)$ \\ $g(X) = \text{sgn}(h(x))$ \\ $g : X \to \{-1 ; 1\}$}

\marge{\couleur{rose}{Hyperplan} \\ $h(x) = 0$}

\begin{itemize}
\item L'hyperplan, c'est un sous-espace de dimension $\mathbb{R}^{p-1}$ (données $X \in \mathbb{R}^p$) ayant pour condit° :
\begin{encadre}
\[ H = \{ x \in \mathbb{R}^p : \langle w, x \rangle + b \} \quad \text{... en gros } h(x) = 0, \forall x \in X \]
\end{encadre}

\couleur{rouge}{\# $\left\{ \begin{array}{l} \text{Sur } \mathbb{R}^1 \text{, H c'est un point} \\ \text{Sur } \mathbb{R}^2 \text{, H c'est une droite} \\ \text{Sur } \mathbb{R}^3 \text{, H c'est un plan} \end{array} \right.$}

\item Pour voir les mauvaises classificat° on peut calculer :
\begin{encadre}
\[ g(x_i) \cdot y_i \geq 0 \]
\end{encadre}

\couleur{rouge}{$g(x_i) = -1$ et $y_i = -1 \Rightarrow g(x)y_i > 0$ \\
$g(x_i) = -1$ et $y_i = 1 \Rightarrow g(x)y_i < 0$ \# \entoure{PB} !}

\couleur{vert}{$\hookrightarrow$ on peut montrer la condit° à minimiser :
\[ \Pr(g(X) \neq Y) = \Pr(g(X)Y = -1) = \Pr(h(X)Y \leq 0) \]}

\item Un échantillon $\{(x_1, y_1), \dots, (x_n, y_n)\}$ est \souligne{linéairement séparable} si
\begin{encadre}
$\exists$ classifieur linéaire (g) qui classe correctement l'échantillon
\end{encadre}

\couleur{rouge}{\entoure{$\forall$} les observat°, c'est \souligne{bien classé}}

\couleur{rouge}{$\hookrightarrow$ C'est quoi le meilleur classifieur optimal ? \\ Il peut exister plusieurs $(w, b)$ qui classe correctement}
\end{itemize}

\dessin[3,81,26,94]{nuage de croix en haut à gauche et de ronds en bas à droite, séparés par 3 droites différentes}

\couleur{vert}{$\hookrightarrow$ C'est quoi le meilleur classifieur entre les 3 ?}
\end{document}
